% ====================================================================
\section{Tutorial — Minimal Polynomial and Diagonalizability}
% ====================================================================

% --------------------------------------------------------------------
\subsection{Recap: Key Facts on Minimal Polynomials}
% --------------------------------------------------------------------

\begin{tipbox}[Recap --- Minimal Polynomial Essentials]
Given $A \in M_{n}(k)$, the \textbf{minimal polynomial} $m_A(x)$ is the monic
polynomial of least degree that annihilates $A$.

\medskip
\textbf{Key results:}
\begin{enumerate}[leftmargin=1.5em,itemsep=3pt]
  \item \textbf{Cayley--Hamilton:} $\chi_A(A)=0$, so $m_A(x)\mid\chi_A(x)$.
  \item \textbf{Common roots:} $m_A$ and $\chi_A$ share exactly the same
        roots (in $\bar{k}$).
  \item \textbf{Stability under field extension:} $m_A(x)$ does not change
        when the coefficient field is enlarged; one may compute and discuss
        $m_A$ over $\mathbb{C}$ even when $A\in M_n(\mathbb{R})$.
  \item \textbf{Diagonalisability criterion:} $A$ is diagonalisable over
        $k$ if and only if $m_A(x)$ has all its roots in $k$ and
        every root is a \emph{simple} root (multiplicity $1$ in $m_A$).
\end{enumerate}
\end{tipbox}

% --------------------------------------------------------------------
\subsection{Worked Examples}
% --------------------------------------------------------------------

\begin{exbox}[Exercise 1 --- Minimal Polynomial of a $3\times 3$ Matrix]
Compute the minimal polynomial of
\[
  A = \begin{pmatrix} 0 & -1 & 0 \\ 1 & -1 & 0 \\ 0 & 0 & 1 \end{pmatrix}.
\]

\textbf{Solution.} The characteristic polynomial is
\[
  \chi_A(x)
  = \det\begin{pmatrix} x & 1 & 0 \\ -1 & x+1 & 0 \\ 0 & 0 & x-1 \end{pmatrix}
  = (x-1)\bigl[x(x+1)+1\bigr]
  = (x-1)(x^2+x+1)
  = x^3-1.
\]
Since $m_A \mid \chi_A$ and they share the same roots, $m_A$ must contain every
irreducible factor of $\chi_A$.  Here $(x-1)$ and $(x^2+x+1)$ are coprime
irreducibles, so
\[
  m_A(x) = (x-1)(x^2+x+1) = \chi_A(x).
\]
\end{exbox}

\begin{exbox}[Exercise 2 --- Diagonalisability over $\mathbb{R}$ and $\mathbb{C}$]
Let $A \in \mathbb{R}^{3\times 3}$ satisfy $\det(A) = -2$ and
\[
  A^4 + 2A^3 + A^2 - 2A - 2I = 0. \tag{$*$}
\]
Decide whether $A$ is diagonalisable over $\mathbb{R}$ and over $\mathbb{C}$.

\medskip
\textbf{Solution.}  Set $P(x) = x^4+2x^3+x^2-2x-2$.
Equation~$(*)$ says $P(A)=0$, so by definition of the minimal polynomial:
\[
  m_A(x) \mid P(x).
\]

\textbf{Step 1: Factor $P(x)$.}  Testing $x=\pm 1$:
\[
  P(1) = 0, \qquad P(-1) = 0.
\]
So $(x-1)(x+1)\mid P$.  Polynomial division gives:
\[
  P(x) = (x-1)(x+1)(x^2+2x+2),
\]
where $x^2+2x+2$ is irreducible over $\mathbb{R}$ (discriminant $= -4 < 0$).

\medskip
\textbf{Step 2: Identify $\chi_A(x)$.}  Since $m_A \mid P$, every eigenvalue
of $A$ is a root of $P$.  Hence the eigenvalues of $A$ lie in
\[
  \{1,\,-1,\,-1+i,\,-1-i\}.
\]
Since $A\in\mathbb{R}^{3\times 3}$, its characteristic polynomial $\chi_A\in\mathbb{R}[x]$
has degree~$3$, and any non-real roots must appear in conjugate pairs.
The only degree-$3$ monic real polynomials with all roots from the above
set are $(x-1)(x^2+2x+2)$ and $(x+1)(x^2+2x+2)$.  Using
$\det A = \lambda_1\lambda_2\lambda_3 = -2$:
\begin{itemize}[leftmargin=1.5em,itemsep=2pt]
  \item $(x-1)(x^2+2x+2)$: product of roots $= 1\cdot(-1+i)(-1-i) = 2
        \neq -2$. \quad\textbf{Rejected.}
  \item $(x+1)(x^2+2x+2)$: product of roots $= (-1)\cdot 2 = -2$.
        \quad$\checkmark$
\end{itemize}
Hence $\chi_A(x) = (x+1)(x^2+2x+2)$.

\medskip
\textbf{Step 3: Identify $m_A(x)$.}  By Theorem~3, $m_A$ and $\chi_A$ share
the same roots: $\{-1,\,-1\pm i\}$.  Over $\mathbb{R}$ the irreducible factors of
$\chi_A$ are $(x+1)$ and $(x^2+2x+2)$, both of which must appear in $m_A$.
Since $m_A\mid\chi_A$, we conclude:
\[
  m_A(x) = (x+1)(x^2+2x+2) = \chi_A(x).
\]

\medskip
\textbf{Conclusion.}
\begin{itemize}[leftmargin=1.5em,itemsep=3pt]
  \item \textbf{Over $\mathbb{R}$:}  $x^2+2x+2$ has no real roots, so not all roots
        of $m_A$ lie in $\mathbb{R}$.  Hence $A$ is \textbf{not diagonalisable
        over $\mathbb{R}$}.
  \item \textbf{Over $\mathbb{C}$:}  The roots $-1,\,-1\pm i$ are all distinct and
        simple.  Hence $A$ \textbf{is diagonalisable over $\mathbb{C}$}.
\end{itemize}
\end{exbox}

\begin{warnbox}[Warning --- Two Easily Confused Facts about Real Polynomials]
\begin{itemize}[leftmargin=1.5em,itemsep=4pt]
  \item A real polynomial of \textbf{odd degree} always has at least one
        \textbf{real} root (by the intermediate value theorem) --- \emph{not}
        a non-real root.  In Exercise~2, the fact that $\chi_A$ has
        non-real roots follows from the eigenvalue constraint and
        $\det A=-2$, \emph{not} from its odd degree.
  \item $P(A)=0$ gives $m_A\mid P$; it does \textbf{not} give $\chi_A\mid P$.
        The correct chain is
        \[
          m_A\mid P \;\Rightarrow\; \text{eigenvalues}\subseteq\operatorname{roots}(P),
        \]
        which is then combined with $A\in\mathbb{R}^{3\times 3}$ and $\det A=-2$ to
        pin down $\chi_A$.
\end{itemize}
\end{warnbox}

% --------------------------------------------------------------------
\subsection{Diagonalisability of $f$ and $f^2$}
% --------------------------------------------------------------------
\begin{propbox}[Proposition --- Diagonalisability of $f$ and $f^2$]
Let $f:\mathbb{C}^n\to\mathbb{C}^n$ be a $\mathbb{C}$-linear isomorphism.  Then
\[
  f \text{ is diagonalisable}
  \;\iff\;
  f^2 \text{ is diagonalisable}.
\]
\end{propbox}

\begin{proofbox}[Proof --- Diagonalisability of $f$ and $f^2$]
\textbf{($\Rightarrow$):}  Suppose $f$ is diagonalisable.
We show that $f^2$ is also diagonalisable.

Since $f$ is diagonalisable, there exists a basis $\mathcal{B}$ of
eigenvectors of $f$, so $[f]_{\mathcal{B}} = D$ is diagonal.  Then
\[
  [f^2]_{\mathcal{B}} = [f]_{\mathcal{B}}^2 = D^2,
\]
which is again diagonal.  Hence $f^2$ is diagonalisable.

\bigskip
\textbf{($\Leftarrow$):}  Suppose $f^2$ is diagonalisable.
We show that $f$ is also diagonalisable.

Since $f^2$ is diagonalisable over $\mathbb{C}$, its minimal polynomial has all roots
in $\mathbb{C}$ and they are simple:
\[
  m_{f^2}(x) = (x-\lambda_1)(x-\lambda_2)\cdots(x-\lambda_r),
  \qquad \lambda_i\in\mathbb{C} \text{ distinct}.
\]
From $m_{f^2}(f^2)=0$:
\[
  (f^2-\lambda_1\,\mathrm{Id})\circ(f^2-\lambda_2\,\mathrm{Id})\circ\cdots\circ(f^2-\lambda_r\,\mathrm{Id}) = 0.
\]
Define $P(x) = (x^2-\lambda_1)(x^2-\lambda_2)\cdots(x^2-\lambda_r)$.
Substituting $f$:
\[
  P(f) = (f^2-\lambda_1\,\mathrm{Id})\cdots(f^2-\lambda_r\,\mathrm{Id}) = 0,
\]
so $P$ annihilates $f$, hence $m_f(x)\mid P(x)$.

\medskip
Since $f$ is an isomorphism, $0$ is not an eigenvalue of $f$, so $\lambda_i\neq 0$
for all $i$.  Over $\mathbb{C}$, each factor splits as
\[
  x^2-\lambda_i = (x-\omega_i)(x-\omega_i'),
\]
where $\omega_i,\omega_i'$ are the two \emph{distinct} square roots of $\lambda_i$
in $\mathbb{C}$ (distinct since $\lambda_i\neq 0$).  Thus
\[
  P(x) = \prod_{i=1}^r (x-\omega_i)(x-\omega_i').
\]

\textbf{All $2r$ roots are distinct.}
\begin{itemize}[leftmargin=1.5em,itemsep=2pt]
  \item Within each pair: $\omega_i\neq\omega_i'$ because $\lambda_i\neq 0$.
  \item Across pairs: if $\omega_i = \omega_j$ (or $\omega_j'$) for $i\neq j$,
        then $\omega_i^2 = \lambda_j$, i.e.\ $\lambda_i=\lambda_j$, contradicting
        the distinctness of the $\lambda_i$.
\end{itemize}

Hence $P(x)$ has $2r$ distinct simple roots in $\mathbb{C}$.  Since $m_f\mid P$, the
minimal polynomial $m_f(x)$ also has all roots in $\mathbb{C}$ and all simple.
Therefore $f$ is \textbf{diagonalisable over $\mathbb{C}$}. \hfill$\blacksquare$
\end{proofbox}